% Section 13.1: Vector fields and their flows | Part: 2 | Depends on: sec:fields, sec:interior | Status: done | Original lines: 2173-2180
%!TEX root = ../../../main.tex
\subsection{Vector fields and their flows}
\begin{formal}[{flow of a vector field}]
A vector field $X\in\Xfield(M)$ is equivalent to an ODE for curves: $\dot\gamma(t)=X_{\gamma(t)}$, \ie in a chart $\dot x^i=X^i(x)$. By existence and uniqueness of ODE solutions $X$ defines a \emph{flow} $\phi_t$ (defined for small $t$ near each point): $\phi_t(x)$ is the solution at time $t$ with initial condition $x$. Properties: $\phi_0=\id$, $\phi_{t_1}\circ\phi_{t_2}=\phi_{t_1+t_2}$ (where defined), $\phi_t^{-1}=\phi_{-t}$, $\frac{\dd}{\dd t}\phi_t(x)=X_{\phi_t(x)}$, and to first order $\phi_\varepsilon(x)=x+\varepsilon X(x)+O(\varepsilon^2)$.
\end{formal}
\begin{intuition}[{wind blowing a leaf}]
$X$ is a stationary wind map; $\phi_t$ moves every point for time $t$ following the wind. The family $t\mapsto\phi_t$ is a one-parameter group of local diffeomorphisms: $X$ is its \emph{infinitesimal generator}.
\end{intuition}
